\section{Discussion and Limitations}
\label{sec:limitations}
Figure~\ref{fig:mmlu-throughput} summarizes completed mean-token measurements
on the RTX 4090 and Arc A770. The shared-prefix study remains pending.
Performance still depends on model size, candidate count and length, and hardware.

The model comparison evaluates the complete MMLU test set but has one run per
model and backend, including each Arc A770 model, so it provides no repeat-run
variance estimate. Model versions,
quantization, placement, and timing boundaries differ across studies, so their
results cannot be treated as a controlled hardware comparison. The planned
prefix study changes both state reuse and cache management; it will not isolate their
individual effects. Shared benchmark coverage also does not make JET and Jev
directly comparable: prompts, scoring, aggregation, and local versus API timing
may differ. A matched constrained-generation baseline is needed to separate the
inference method from the underlying model's capability.
Comparison with a fixed Jev version should use the same contexts, candidate
meanings, answer types, and evaluation rules. Performance measurements should
report warm p50/p95 latency, throughput at fixed concurrency, startup cost,
and failure rates, distinguishing local inference from remote service latency.

Mean token log-probability reduces the cumulative penalty on longer answers,
but remains sensitive to wording, tokenization, and answer priors. Easily
predicted padding can increase a candidate's mean score; quotation marks and
escaped content also contribute to the scored token count. Length normalization
therefore does not guarantee length-invariant decisions. The resulting softmax
weights are not candidate-restricted sequence likelihoods, and their
normalization does not establish calibration~\cite{calibration}; this requires
separate empirical evaluation for categorical, binary, and ordinal outputs.
NLL is a diagnostic, not a substitute for
calibration curves, Brier scores, or selective-prediction evaluation.
Pretraining data may also overlap benchmark content.
The fine-tuning experiment excludes the frozen MMLU requests and exact
normalized overlaps, but cannot rule out paraphrases or pretraining
contamination. Its public text tasks do not establish improved visual,
ordinal, or closed-loop decision quality. The training mixture, adapter rank,
and quantization have no ablation study.

Sampler-free selection removes sampling variance, but numerical differences
across hardware, quantization, kernels, and batch layouts can remain. Exact
agreement applies only to the tested configurations. Hybrid continuations are
currently evaluated serially; concurrency across independent questions requires
state isolation and numerical validation. Multimodal accuracy and closed-loop
control remain unmeasured. Raw experiment artifacts are retained outside version
control, so independent replication requires those records or regeneration under
the documented conditions.
